\subsection{All sections over the split-prime source at seven}

\begin{statement}{E.98}{Proved}{The complete seven-and-three class-two locus and the full-to-metabelian ramified fibres}
We retain the actual based geometric pro-three group $N$ of
the cubic with its nine cusps removed, the marking of E.95,
and its genuine arithmetic action. Put
\[
 N_2=N/\gamma_3N,\qquad N_{\mathrm m}=N/N'',\qquad
 V=N^{\mathrm{ab}},\quad M=\bigwedge^2V,\quad
 W=\langle e_d:d\ne O\rangle\subset V.
 \tag{E98.1}
\]
Here $V$ has rank ten and $W$ has rank eight over $\mathbb Z_3$.
For each of these three target groups $H$, let
$\operatorname{Sect}_{37}(H)$ be the geometric conjugacy
classes of continuous global sections whose cocycles are
unramified outside $\{3,7\}$. The following coordinates
compute these stated global section spaces.

Write $T=\Phi_\tau$ and $J=\Phi_c$ for the actual
Frobenius and complex-conjugation actions specified below.
With $T_1=T|V$, the horizontal abelian space is
\[
 P_1=V/(I-T_1)V^{J_1}
       \simeq C_3^3\oplus\mathbb Z_3^6.
 \tag{E98.2}
\]
Use the retained E.37 Smith representatives $U\in V$ for
$P_1$. The entire abelian section space is $P_1\times W$.
Set $e_O=-\sum_{d\ne O}e_d$, and let $k_d\in V$ be
the actual E.95 Frobenius cusp columns; set $k_O=0$.
The image of $\operatorname{Sect}_{37}(N_2)$ in this
abelian space consists exactly of the following pairs:
\[
\begin{array}{ll}
 v\in9W:& U\text{ arbitrary};\\[2pt]
 v\equiv r e_D\pmod{9W},\quad1\le r\le8:
   &r(U-k_D)\in (T_1-I)V+\mathbb Z_3e_D .
\end{array}
 \tag{E98.3}
\]
The nine lines have disjoint nonzero residues modulo nine,
so $D,r$ in the second row are unique. Every local quotient
in this formula is given by an explicit Smith map with
both inverse transformations:
\[
 V/\bigl((T_1-I)V+\mathbb Z_3e_D\bigr)
               \simeq C_3^5\oplus C_9^2
                    \qquad(D\in E[3]).
 \tag{E98.4}
\]
In particular the complete relative image and the lattice
admitting every horizontal input are, respectively,
\[
 \bigcup_{D\in E[3]}(\mathbb Z_3e_D+9W),
                  \qquad 9W.
 \tag{E98.5}
\]
The first set has exactly $73$ residues modulo nine.

Over every pair in (E98.3), the complete geometric fibre is
\[
             \mathbb Z_3^{25}\times C_3^9\times C_9^5.
 \tag{E98.6}
\]
Over a fixed $v$ in (E98.5), the complete fibre of the
relative-inertia map is
\[
 \mathbb Z_3^{31}\times C_3^{9+t(v)}\times C_9^5,
 \qquad
 t(v)=\begin{cases}3&v\in3W,\\0&v\notin3W.\end{cases}
 \tag{E98.7}
\]
These are chosen continuous coordinates with explicit
inverses. The proof gives the genuine conjugation changes
between representatives and all origins; it does not
assign an abelian group operation to the nonabelian
section space.

For the full and intrinsic metabelian targets, the
zero-inertia subspace is the complete unramified space
computed in E.41. Every ramified class has exactly the
coordinates
\[
 \coprod_{D\in E[3]}
  (\mathbb Z_3\setminus\{0\})\times\{0,1,2\}
                                  \times H^J\backslash H,
             \qquad H=N,\ N_{\mathrm m}.
 \tag{E98.8}
\]
The actual forward map uses the chosen positive peripheral
element $\delta_D$ and its tangential Frobenius $z_D\tau$:
\[
 \begin{split}
 a(t)&=g\delta_D^r g^{-1},\\
 a(\tau)&=g\delta_D^b z_DT(g)^{-1},\qquad a(c)=1,\\
 a(t')&=J(a(t))^{-1},\qquad a(\tau')=J(a(\tau)).
 \end{split}
 \tag{E98.9}
\]
The path and peripheral elements for $N_{\mathrm m}$ are
the actual projections of those for $N$.

There are explicit coordinates for the last factor:
\[
 \begin{split}
 H^J\backslash H&\longrightarrow
       \operatorname{Anti}_J(H):=\{x\in H:J(x)=x^{-1}\},\\
 [g]&\longmapsto g^{-1}J(g),\qquad
 x\longmapsto[x^{-1/2}].
 \end{split}
 \tag{E98.10}
\]
Moreover these spaces have recursive Hall and Chen
coordinates
\[
 \operatorname{Anti}_J(N)\simeq
       \prod_{m\ge1}\mathbb Z_3^{\nu_m},\qquad
 \operatorname{Anti}_J(N_{\mathrm m})\simeq
       \prod_{m\ge1}\mathbb Z_3^{\bar\nu_m}.
 \tag{E98.11}
\]
All ranks are evaluated by the following formulas. Put
$c_m=[z^m](1-z)^{-4}(1+z)^{-6}$ and
$a_d=-2$ for odd $d$, $a_d=10$ for even $d$. Then
\[
 \begin{split}
 \nu_m&=\frac1{2m}\sum_{d\mid m}\mu(d)
                       \bigl(10^{m/d}-a_d^{m/d}\bigr),\\
 \bar\nu_1&=6,\qquad
 \bar\nu_m=\frac12\left((m-1)\binom{m+8}{m}
                                  +2c_{m-1}+c_m\right)
                                    \quad(m\ge2).
 \end{split}
 \tag{E98.12}
\]
These coordinates retain the action on each actual
nilpotent quotient, including its higher commutator terms.

The map from full to metabelian ramified sections is
\[
       (D,r,b,[g])\longmapsto(D,r,b,[\bar g]).
 \tag{E98.13}
\]
It is onto, and every fibre, with its explicit origin, is
\[
 (N'')^J\backslash N''\simeq
       \prod_{m\ge4}\mathbb Z_3^{\nu_m-\bar\nu_m}.
 \tag{E98.14}
\]
The degree-four factor has rank $504$. Thus (E98.14)
evaluates the full global fibre of the displayed map.

There are also actual global class-two sections in this
source with inertia $9(e_{A_0}+e_{B_0})$ and horizontal
coordinate zero. Every such section has no intrinsic
metabelian lift and no full lift. The explicit central
inertia and the complete tame relation check appear below.

\emph{Manuscript: \S51.16, folios733v--736v, scans172--178,
formulas(371)--(379), Part E(c),(d) and the fourth exactness
follow-up. The unrestricted global and higher-stage
requirements remain recorded in E.23.}
\end{statement}
\begin{proof}
\emph{The actual arithmetic source and both comparison maps.}
Let $K=\mathbb Q(\zeta_3)$, let $\mathfrak p,\bar{\mathfrak p}$
be the primes over seven, and put $\pi=1+3\zeta_3$.
We use the actual normalized Frobenii of E.56 at these
two prolongations. Its source isomorphism, specialized
to seven, is
\[
 G_{37}=\left(D_7\mathop{\ast}_3D_7'\right)
                                    \rtimes\langle c\rangle,
 \quad
 \begin{cases}
 D_7=\langle t,\tau\mid\tau t\tau^{-1}=t^7\rangle_3,\\
 D_7'=\langle t',\tau'\mid\tau't'(\tau')^{-1}=(t')^7\rangle_3,\\
 c\tau c=\tau',\qquad ctc=(t')^{-1}.
 \end{cases}
 \tag{E98.15}
\]
Its characters are $\chi(\tau)=\chi(\tau')=7$,
$\kappa_3(3)(\tau)=1$, $\kappa_3(3)(\tau')=-1$,
and $\kappa_3(7)(t)=\kappa_3(7)(t')=1$.
The latter cocycle vanishes at both chosen Frobenii.
The two former Frattini characters have determinant
$-2(7-1)/3$, a three-adic unit. Their actual images
therefore form a free basis of the group unramified
outside three computed in E.34.

We recall the exact kernel comparison used in E.56.
\cite[Theorem2.4, its proof and the remark after
Proposition2.1]{E_WingbergRiemann} identifies the kernel
with the continuous free product of the actual inertia
groups over $P/\langle\tau\rangle$ and
$P/\langle\tau'\rangle$. The two bundle relations are
$(p\tau^z,t^u)=(p,t^{7^zu})$ and their primed versions.
The coset tree for $D_7\ast_3D_7'$ has these same
stabilizers and these same conjugation maps. Maps
induced by the four actual generator inclusions give
inverse maps on these bundles and the identity on
the quotient $P$. This proves (E98.15) with its specified
local subgroups. All primes above three are in the
ramification set, and the cyclotomic tower has nontrivial
unramified completion at both added primes, as required
in that theorem. The actual Frobenii and inertias also
belong to the full relative arithmetic group; the same
relations give a reverse map whose composite with the
arithmetic projection is identity. It commutes with $c$.

The actual geometric action is unramified outside three
by E.34 and the based tangent comparison of E.78. Its
image over $K$ is pro-three. Pairing a cocycle with this
action gives a homomorphism into a pro-three semidirect
product. Consequently a cocycle in the stated
ramification locus factors through (E98.15); inflation
gives the inverse construction on the original
arithmetic source. Both maps preserve geometric
conjugacy. Thus relations in (E98.15) construct the
requested global sections in this locus.

An odd profinite group has unique continuous square
roots: on each finite quotient of odd exponent $e$,
the inverse of squaring sends $x$ to $x^k$ for
$2k\equiv1\pmod e$. These inverses commute with
quotient maps. A cocycle value $z=a(c)$ satisfies
$J(z)=z^{-1}$. Geometric conjugation by $z^{-1/2}$
makes this value one. The residual conjugators are
exactly $H^J$. Set $u=a(\tau)$ and $a=a(t)$.
The other two values are forced as in (E98.9), and
the first tame relation becomes
\[
                 uT(a)u^{-1}=a^7.
 \tag{E98.16}
\]
The second is its $J$-image. Hence normalized sections
are exactly the solutions of (E98.16), modulo
$(u,a)\mapsto(huT(h)^{-1},hah^{-1})$, $h\in H^J$.
The free-product universal property proves sufficiency
and continuity, in addition to necessity.

\emph{The exact class-two differential.}
In the E.95 basis $(f_1,f_2,A_0,A_1,A_2,B_0,B_1,B_2,C_1,C_2)$,
put $R=\left(\begin{smallmatrix}0&-1\\1&-1\end{smallmatrix}\right)$
and retain the gluing matrix $U_{\rm gl}$ of E.37.
The actual abelian Frobenius is
\[
 T_1=\begin{pmatrix}I_2+3R&0\\
        U_{\rm gl}R-2U_{\rm gl}&7I_8\end{pmatrix}.
 \tag{E98.17}
\]
Use the actual $J$-linear class-two coordinates of E.95,
with product $(v,m)(v',m')=(v+v',m+m'+[v,v']/2)$.
Write $B=\bigwedge^2T_1$ and $J_2=\bigwedge^2J_1$.
For the actual cusp path columns,
$T(e_d,0)=(7e_d,-7[k_d,e_d])$.
The abelian part of (E98.16) forces $v\in W$;
the proper block $T_1-7I$ is rationally invertible.
Its central part is exactly
\[
 \mathcal d m+\Delta(U,v)=0,\qquad
 \mathcal d=B/7-I,
 \quad\Delta(U,v)=\sum_{d\ne O}v_d[U-k_d,e_d].
 \tag{E98.18}
\]
Here $u=(U,z)$ and $a=(v,m)$; the central $z$ cancels.
Only the actual action on peripheral inertia occurs
in (E98.18), and those columns are fully evaluated in E.95.

Let $j=(3I_2;U_{\rm gl})$ and $w=j(f_1)\wedge j(f_2)$.
It is primitive, $Bw=7w$, and $J_2w=-w$.
Order the forty-five wedge coordinates by increasing
pairs of the displayed ten indices. The coefficient
of $w$ at $A_0\wedge A_2$ is one; call this position $p$.
Let $Q$ subtract the $p$-th coordinate times $w$ and
then delete that coordinate, and let $S$ insert zero
at that position. Let $P$ delete the $f_1\wedge f_2$
coordinate. The actual maps satisfy
\[
 QS=I_{44},\quad SQ=I_{45}-w e_p^{\mathsf t},\quad
 Qw=0,\qquad D=P\mathcal dS.
 \tag{E98.19}
\]
The image of $\mathcal d$ is in
$M_0=\ker(e_{f_1\wedge f_2}^{\mathsf t})$.
The retained exact Smith identity $LDR=D_0$ has
\[
 D_0=\operatorname{diag}(1^{[14]},3^{[18]},9^{[12]}),
 \qquad D^{-1}=R D_0^{-1}L.
 \tag{E98.20}
\]
The certificate checks both inverse products and all
four change-of-basis inverse products. Thus $\mathcal d$
has precisely kernel $\mathbb Z_3w$, and a particular
solution is
\[
                       m_0(U,v)=-SD^{-1}P\Delta(U,v).
 \tag{E98.21}
\]
It is integral exactly when the eighteen modulo-three
and twelve modulo-nine Smith congruences vanish. Every
integral solution is uniquely $m_0(U,v)+\lambda w$,
$\lambda\in\mathbb Z_3$; its inverse scalar is its
$p$-th coordinate. The full cokernel obstruction in
this equation is $C_3^{18}\oplus C_9^{12}$.

\emph{Every residue and the nine exact cusp lattices.}
The E.37 first-degree Smith map for
$(I-T_1)(I+J_1)$ has nonzero diagonal $1,3,3,3$.
Its inverse target matrix, with the unit column
removed, gives the ten-by-nine representative matrix
$R_1$ used here. Its first three coordinates are reduced
modulo three and its last six are arbitrary in $\mathbb Z_3$.
For each of the eight $d$, the certificate records the
thirty-by-ten matrix of the nonunit rows of
\[
       LP[\,\cdot\,,e_d],[-k_d\mid R_1].
 \tag{E98.22}
\]
All gauge columns vanish in their exact cokernel, and
three times each finite horizontal column vanishes.
This proves the descent to $P_1\times W$, with its
finite carries included.

All coefficients of (E98.21) have three-adic
denominator dividing nine. Therefore its integrality
depends on $v\bmod9$ and on the indicated horizontal
residues. The retained two-digit elimination tests
all $6561$ first digits. Exactly nineteen admit a
horizontal value. For each of those it tests all
$6561$ second digits, retaining the divided constant
and divided-kernel carry, and solves the remaining
linear system over $\mathbb F_3$. The zero first digit
has nineteen admissible second digits; each of the
eighteen nonzero first digits has three. All seventy-three
resulting horizontal witnesses are substituted into
every original congruence and into the full rational
inverse (E98.21). Their inertia vectors are exactly
\[
          \{r e_D\bmod9W:D\in E[3],\ 0\le r<9\}.
 \tag{E98.23}
\]
The set equality is checked directly, with all actual
vectors retained. Thus the finite enumeration proves
exclusion as well as construction of the relative image.

Here is the exact horizontal criterion on these residues.
For every $D$, the projection of $w$ in
$\bigwedge^2(V/\mathbb Z_3e_D)$ is primitive. The
certificate supplies an integral unimodular basis
containing $e_D$, its inverse, and a unit wedge
coefficient after deleting that direction. Its
ten-by-eleven matrix $[T_1-I\mid e_D]$ has Smith
diagonal $1^{[3]},3^{[5]},9^{[2]}$, with an explicit
rational right inverse. This proves (E98.4) as an
isomorphism with specified maps.

The actual cusp origin for $D\ne O$ is $U=k_D$,
and for $O$ it is $U=0$ with its retained central
peripheral coordinate. In either case (E98.21) at
the origin is integral. Put
$h=(T_1-I)^{-1}(U-k_D)$ over $\mathbb Q_3$; this
inverse exists because none of the eigenvalues of
$T_1$ is one. The identity
\[
           \mathcal d[h,e_D]=[(T_1-I)h,e_D]
 \tag{E98.24}
\]
gives the change of particular solution as
$-SQ[rh,e_D]$. Integrality of this expression is
equivalent to $rh\in V+\mathbb Q_3e_D$.
Indeed a rational multiple of $w$ making the wedge
integral projects to a multiple of the primitive
projected $w$; that multiple is integral. Removing
it shows $[rh,e_D]\in\bigwedge^2V$, and the integral
basis containing $e_D$ identifies this with the
stated condition on $rh$. Applying $T_1-I$, and using
$(T_1-I)e_D=6e_D$, identifies it exactly with
$r(U-k_D)\in(T_1-I)V+\mathbb Z_3e_D$.
Here $V\cap\mathbb Q_3e_D=\mathbb Z_3e_D$ by
primitivity. Adding $9W$ to the inertia changes
(E98.21) integrally for every horizontal value.
This proves both rows of (E98.3).

The map $P_1\to V/((T_1-I)V+\mathbb Z_3e_D)$ is
onto, since it is induced by the identity on $V$.
For a unit $r$ its kernel has index $3^9$, and for
$r\in3\mathbb Z_3\setminus9\mathbb Z_3$ the
kernel of multiplication by $r$ after that map has
index nine. The certificate evaluates all twenty-seven
torsion inputs of $P_1$: only zero survives in the
first kernel, and all survive in the second.
Thus their additive types are respectively
$\mathbb Z_3^6$ and $\mathbb Z_3^6\oplus C_3^3$;
for inertia in $9W$ the whole $P_1$ occurs. Exact
row-kernel intersection of every constant and
horizontal coefficient in (E98.22) has Hermite
basis $9I_8$. This independently proves the second
claim of (E98.5). The quotient (E98.4) also shows
why every nonzero inertia residue imposes a
nonzero horizontal condition.

\emph{All class-two fibres and actual inverse coordinates.}
For every allowed $(U,v)$, choose $m$ from (E98.21)
and an arbitrary central $z$. After the horizontal
representative is fixed, residual conjugators are
exactly $M^{J_2}$, and they act by
$z\mapsto z+(I-B)k$. Use the integral $J_1$ eigenspace
basis of E.95, of ranks four and six, and its wedge
splitting. The twenty-one same-sign wedges give
the exact basis matrix $M_+$ of $M^{J_2}$.
For $C=(I-B)M_+$, the full Smith data are
\[
 L_C C R_C=D_C,\qquad
 D_C=\operatorname{diag}(1^{[7]},3^{[9]},9^{[5]}),
                    \quad C\text{ has size }45\times21.
 \tag{E98.25}
\]
All four inverse products and forty complete
coordinate-and-conjugator round trips are retained.
Use zero at the seven unit entries, standard
residues at the fourteen finite entries, and
arbitrary three-adic values at the last twenty-four
entries, then multiply by $L_C^{-1}$. The inverse
multiplies by $L_C$ and reduces the first twenty-one
entries. Its correcting conjugator is
$R_C D_C^{-1}$ times that difference, in those
twenty-one coordinates. Together with the scalar
$\lambda$ of (E98.21), this proves (E98.6).

For completeness the first-degree change also
has a literal class-two inverse. The unique
$h\in V^{J_1}$ carrying an arbitrary $U$ to its
Smith representative is obtained from the
first-degree Smith inverse. Conjugate by $(h,0)$,
using the genuine action $T(h,0)=(T_1h,\eta_T(h))$.
The central changes are exactly
\[
 \begin{split}
 z&\longmapsto z-\eta_T(h)
       +\tfrac12[h,U]-\tfrac12[h+U,T_1h],\\
 m&\longmapsto m+[h,v].
 \end{split}
 \tag{E98.26}
\]
The handle-action term is the actual arithmetic
conjugation of the chosen lift. Formula (E98.26)
specifies its use and preserves it. The inverse
does not replace that action by its graded matrix.
After this change apply (E98.21) and (E98.25).
Conversely these formulas construct every pair
satisfying the actual tame relation and hence
every global section in this locus. Uniqueness
of $h$ follows from rational injectivity of
$(I-T_1)|V^{J_1}$, and uniqueness of the residual
central conjugator follows from the full column
rank in (E98.25). Thus no stabilizer quotient is
omitted. Combining the horizontal kernels already
computed with this central fibre gives (E98.7).
All maps are continuous on the indicated integral
lattices and their finitely many residue strata.

For a concrete nonlifting section, put
\[
 U=0,\quad v=9(e_{A_0}+e_{B_0}),\quad
 u=1,\quad a(t)=(v,m_0(0,v)),\quad a(c)=1,
 \tag{E98.27}
\]
and use (E98.9) for the primed values. The full
forty-five-coordinate vector (E98.21) is integral
and is retained in the certificate; direct
substitution gives $\mathcal dm_0=-\Delta(0,v)$.
This proves the tame word exactly, including all
central coordinates. Every one of the nine matrices
$[v\mid e_D]$ has rank two. E.61 and E.66 classify
the full and intrinsic metabelian local tame
sections and put every nonzero abelian inertia
on exactly one of these cusp lines. The actual
projections $N\to N_{\mathrm m}\to N_2$ preserve
these inertia vectors. Restricting a global lift
to the chosen place over seven would give such
a local lift, which the computed vector excludes.
The same argument applies to every central class
over this abelian input.

\emph{The full and metabelian ramified packets.}
For either $H=N$ or $H=N_{\mathrm m}$, E.61 and
E.66 apply to every nonidentity inertia in
(E98.16). In the metabelian case their proof
first excludes nonidentity inertia with zero
abelianization and then treats all nonunit
exponents. They give
$a=g\delta_D^r g^{-1}$, $r\ne0$, and
$C_H(\delta_D^r)=\langle\delta_D\rangle$.
The actual cusp path satisfies
$z_DT(\delta_D)z_D^{-1}=\delta_D^7$.
Substitution in (E98.16) and this centralizer
equality give exactly its second value in
(E98.9), with $b\in\mathbb Z_3$.
The abelian inertia determines $D$ and $r$.

Two representatives for fixed $D,r$ have the
same inertia precisely when $g'=g\delta_D^k$.
Using the displayed actual Frobenius relation
in the second value shows that their values
of $u$ coincide precisely when $b'=b+6k$.
Choose the unique $b\in\{0,1,2\}$ representing
its class modulo $6\mathbb Z_3$. This determines
$k$ uniquely, so removes the whole right-cusp
ambiguity. The residual global conjugations
are exactly left multiplication by $H^J$.
This proves (E98.8), with forward map (E98.9)
and inverse obtained from these unique data.
The $a=1$ part of (E98.16) is exactly E.41's
unramified section quotient. These two parts
exhaust the stated full and metabelian spaces.

The map (E98.10) is left-$H^J$ invariant.
Equality of the images of $g,g'$ gives
$J(g'g^{-1})=g'g^{-1}$, proving injectivity.
For an anti element $x$, its inverse image
$x^{-1/2}$ satisfies $g^{-1}J(g)=x$.
This proves surjectivity and both inverse
formulas; continuity follows from that of
the unique roots. On every shell
$v_3(r)=n$, the parameter space (E98.8) is
compact and its target geometric orbit
space is Hausdorff, so the continuous
bijection is a homeomorphism. The shells
are clopen in the ramified subspace, as
detected by abelian inertia. Their union
therefore has the stated topology.

\emph{The exact full-to-metabelian global fibre.}
Put $R_0=N''$. Formula (E98.13) follows by
applying the actual quotient to every word
in (E98.9). It is onto. We also have a
literal lift of every fixed element:
for $\bar h\in N_{\mathrm m}^J$, choose
$h_0\in N$, put $k=h_0^{-1}J(h_0)\in R_0$,
and set $h=h_0k^{1/2}$. Since $J(k)=k^{-1}$,
$J(h)=h$ and $h$ projects to $\bar h$.
Choose a lift $g_0$ of a target representative
$\bar g$. Every source coset over it has
a representative $kg_0$, $k\in R_0$:
lift and remove its left fixed factor by
the preceding formula. Two such representatives
are equivalent exactly when $k'=hk$ with
$h\in R_0^J$. The actual fibre map is thus
\[
              R_0^J\backslash R_0\longrightarrow
                 \operatorname{fib}_{[\bar g]},
                      \qquad[k]\longmapsto[kg_0].
 \tag{E98.28}
\]
Its inverse is the described removal of the
fixed factor. It is a continuous bijection
between compact Hausdorff spaces, hence a
homeomorphism.

\emph{All-degree coordinates for the remaining explicit quotients.}
We give the construction underlying (E98.11)--(E98.14).
The integral lower-central comparison for $N$ is
the free Lie algebra on $V$, as proved in E.40--E.41.
For $N_{\mathrm m}$ we use the exact completed Fox
map of (E41.13),
\[
 N'/N''\longrightarrow
   \ker\bigl(\mathbb Z_3[[X_1,\ldots,X_{10}]]^{10}
                       \xrightarrow{(X_i)}
                    \mathbb Z_3[[X_1,\ldots,X_{10}]]\bigr).
 \tag{E98.29}
\]
The homogeneous filtration equality proved in
(E41.15) identifies its degree $m\ge2$ quotient
with
$\ker(V\otimes\operatorname{Sym}^{m-1}V
                     \to\operatorname{Sym}^{m}V)$.
That proof gives the actual commutator-to-syzygy
map, its integral Chen basis and its induced
pro-three topology. Thus all these lattices
are free, with no extra completion torsion.
For $R_0$ use the induced filtration
$R_0\cap\gamma_mN$. Its graded quotient is
exactly the kernel of the two preceding
graded maps: an element with zero graded
image has image in $\gamma_{m+1}N_{\mathrm m}$;
lift that image in $\gamma_{m+1}N$ and divide
to obtain an element of $R_0$ with the same
graded coordinate. This also proves the
map direction needed for surjectivity onto
that kernel. It is saturated because the
target graded lattice is free.

Here is an explicit coordinate procedure
for each of the three resulting groups.
At one central filtration step, start
with an anti element in the lower quotient.
Use the ordered Hall or Chen lifts, with
zero new coordinates, to obtain a lift $u$.
For $R_0$ use the saturated kernel basis
of the actual Hall-to-Chen matrix and its
lifted commutators. The element
$d=uJ(u)$ belongs to the central new
lattice. Centrality gives $J(d)=d$.
Consequently
\[
                  u d^{-1/2}
 \tag{E98.30}
\]
is an anti lift. All other anti lifts are
uniquely its products by the negative
eigenspace of that central lattice.
The projectors $(I+J_m)/2$ and $(I-J_m)/2$
are integral, so their exact lattice
bases and inverses give the new coordinates.
For the inverse procedure, subtract the
previous canonical lift and read the
negative central coordinates. This proves
the two mutually inverse recursive maps,
including every nonlinear carry.

Ordered commutator power maps are continuous
for three-adic exponents; the integral
collection identities give continuous Hall
and Chen coordinate sections at each finite
nilpotent stage. The corresponding kernel
bases supply the same sections for $R_0$.
Every finite three-group is nilpotent.
Hence the intersection of the closed
lower-central terms of these groups is
trivial, and compactness makes the map
onto the inverse limit of these quotients
surjective. For $R_0$ the induced filtration
has the same properties by its inclusion
in $N$. Passing to that inverse limit in
(E98.30) proves the product homeomorphisms.

It remains to evaluate their ranks. On $V$,
$J_1$ has positive rank four and negative
rank six. Integral Hall independence and
the tensor-algebra PBW identity give
\[
 \operatorname{rank}\operatorname{Lie}_m(V)
     =\frac1m\sum_{d\mid m}\mu(d)10^{m/d},\qquad
 \operatorname{tr}(J_m)
     =\frac1m\sum_{d\mid m}\mu(d)a_d^{m/d}.
 \tag{E98.31}
\]
For example the second identity follows by
taking the trace generating series in PBW,
expanding the formal logarithm, and applying
M\"obius inversion; $\operatorname{tr}(J_1^d)=a_d$.
In the metabelian syzygy sequence the
multiplication map onto $\operatorname{Sym}^mV$
is onto over $\mathbb Z_3$, since each
monomial is the product of a variable and
a monomial of degree $m-1$. Its kernel
therefore has rank $(m-1)\binom{m+8}{m}$
and trace $-2c_{m-1}-c_m$.
Taking one half of rank minus trace gives
(E98.12). Exactness of the saturated
kernel sequence gives negative rank
$\nu_m-\bar\nu_m$ for $R_0$.
The derived ideal has its first nonzero
second-derived component in degree four;
the same formulas give zero in degrees
one through three and $1248-744=504$
in degree four. The program evaluates
the fourteen first exact character rows
and their integrality. Applying (E98.10)
to $R_0$, followed by (E98.30), proves
the final product formula (E98.14).
\end{proof}
